
%%% Local Variables:
%%% mode: latex
%%% TeX-master: "../main"
%%% End:

\begin{ack}

光阴荏苒，七载求学路已至终点。回首在华中科技大学度过的两千余个日夜，喻家山的四季轮转与东湖的微风细雨，皆已化作心中最深切的感激与敬畏。

首先，我要向我的导师夏天教授表达最深切的谢意。 老师严谨的治学态度、敏锐的学术洞察力以及对科研近乎热忱的追求，始终深深影响着我。在遇到科研瓶颈时，是老师的悉心指引让我拨云见日；在初入学术殿堂时，是老师的宽容与鼓励让我坚定了前行的信心

其次，感谢实验室同甘共苦的伙伴们。 感谢蔡训辉、胡力夫、齐宇轩、钟睿芸、姜晨、王夏宇等师兄师姐在科研初期的领航；感谢方一舟、付丛、赵佳飞、曹威等同级挚友的并肩战斗；感谢江宇、梁川、杨毅等师弟带来的蓬勃活力。实验室里无数个灯火通明的夜晚，关于实验数据的激烈讨论与论文逻辑的反复推敲，皆是我学生时代最宝贵的精神财富。

感谢一路同行、见证成长的朋友们。 感谢本科好友杨福康、张睿哲等，以及高中好友崔龙飞、王宇航、郭志浩等。即便身处不同领域，你们的陪伴与支持亦是我抵御孤独的良药。

特别感谢我的女朋友李阳。 感谢你在每一个压力剧增的时刻给予我的温暖支撑，这份温柔与包容，是我在科研长跑中不竭的动力。

最要感谢我的家人。 感谢父母多年来的默默付出与无私支持，你们的爱如静水流深，是我在风雨中勇往直前最坚实的后盾。

最后，感谢在百忙之中审阅本论文的专家教授们。 你们的批评与建议将是我未来学术道路上不断进步的阶梯。

凡是过往，皆为序章。感谢华中科技大学，赋予我独立之人格、奋斗之勇气。

\end{ack}
